\chapter{Cloud Computing Fundamentals and Rationale}
\label{chap:cloud_overview}

\section{Service Models in Cloud Computing}
Cloud computing abstracts raw physical infrastructure into accessible, programmable digital services. The Cloud Security Alliance (CSA) and NIST classify cloud offerings into three primary service models:
\begin{enumerate}
    \item \textbf{Infrastructure as a Service (IaaS)}: Provides raw compute instances, virtual disks, and software-defined networking primitives. Users maintain full control over operating systems, runtime middleware, security patches, and application deployments, while the provider manages physical servers and data center facilities. (e.g., Amazon EC2, AWS EBS).
    \item \textbf{Platform as a Service (PaaS)}: Delivers pre-configured execution environments where developers deploy code without managing OS maintenance or web server binaries. (e.g., AWS Elastic Beanstalk).
    \item \textbf{Software as a Service (SaaS)}: Delivers fully managed, consumer-facing end-user applications hosted entirely in the cloud. (e.g., Google Workspace, Microsoft 365).
\end{enumerate}

In recent years, \textbf{Serverless Computing} (Function-as-a-Service, FaaS) has emerged as an evolution of PaaS. In serverless computing, developers deploy stateless event-driven functions (e.g., AWS Lambda) while the cloud vendor manages automatic scaling, process isolation, OS patching, and execution scheduling.

\section{Core Attributes of Cloud Architecture}
Modern cloud applications leverage key architectural attributes that distinguish them from traditional deployments:
\begin{itemize}
    \item \textbf{Scalability and Elasticity}: Scalability refers to the system's ability to handle growing workloads by increasing capacity. Elasticity is the dynamic adaptation of capacity—scaling out during traffic surges and scaling in when load subsides—ensuring resource allocation strictly tracks demand.
    \item \textbf{High Availability and Fault Tolerance}: Cloud infrastructure is distributed across geographically isolated Availability Zones (AZs). Applications designed with multi-AZ replication survive hardware failures without downtime.
    \item \textbf{Pay-As-You-Go Financial Model}: Replaces upfront capital expenditure (CapEx) with operational expenditure (OpEx), charging consumers only for exact compute duration, request volume, and bytes stored.
    \item \textbf{Shared Responsibility Model}: Defines the division of security obligations between the cloud provider and the consumer. The provider guarantees security \textit{of} the cloud (hardware, facilities, hypervisors), while the consumer manages security \textit{in} the cloud (IAM policies, database encryption, network access rules, application code).
\end{itemize}

\section{Application to Institutional Campus Management}
The College Campus Management System exemplifies an ideal workload for cloud migration:
\begin{table}[htbp]
\centering
\caption{Traditional On-Premises Deployment vs. Cloud-Native AWS Architecture}
\label{tab:onprem_vs_cloud}
\small
\begin{tabular}{L{3.5cm} L{5.5cm} L{6.0cm}}
\hline
\textbf{Architectural Attribute} & \textbf{Traditional On-Premises Setup} & \textbf{Cloud-Native AWS Architecture} \\ \hline
\textbf{Capacity Planning} & Over-provisioned for peak load; idle 85\% of time. & Elastic serverless scaling tracking live HTTP request rate. \\ \hline
\textbf{Capital Cost (CapEx)} & High upfront server procurement and rack costs. & \$0 upfront; operational billing based on usage. \\ \hline
\textbf{High Availability} & Single server failure causes institutional downtime. & Multi-AZ replication for RDS, S3, and Lambda functions. \\ \hline
\textbf{Asset Storage} & Local hard drives prone to disk corruption. & Amazon S3 with 99.999999999\% (11 9s) durability. \\ \hline
\textbf{Security Maintenance} & Manual OS security patching and firewall rules. & Automated AWS managed OS patching and IAM RBAC rules. \\ \hline
\end{tabular}
\end{table}
